% from ega4-17-fr.tex
\oldpage[IV-4]{77}

autre que $Y-f(X-E)$, donc (ii) résulte de (i) et du théorème de Chevalley lorsque $f$ est de présentation finie (1.8.4). De même, si $h:Y\to S$ est le morphisme structural, l'ensemble des $s\in S$ vérifiant la propriété correspondante considérée dans (iii) est $S-h(Y-F)$, donc (iii) découle encore de (ii) et du théorème de Chevalley lorsque $f$ et $h$ sont de présentation finie. Il suffit donc de prouver les assertions de (i).

Prouvons d'abord (i) lorsqu'il s'agit de la propriété d'être \emph{lisse}. La question étant locale sur $X$, on peut se borner au cas où $S=\Spec(A)$, $X=\Spec(B)$ et $Y=\Spec(C)$ sont affines, $B$ et $C$ étant des $A$-algèbres de présentation finie. Raisonnant comme au début de (9.9.1) et utilisant (17.7.2, (ii)), on se ramène au cas où $A$ est noethérien. En vertu de ($0_{\mathrm{III}}$, 9.2.3), on est ramené à voir que si $x\in E$ (resp. si $x\notin E$), il existe un voisinage $V$ de $x$ dans $\overline{\{x\}}$ contenu dans $E$ (resp. dans $X-E$). Désignant par $g:X\to S$ et $h:Y\to S$ les morphismes structuraux, on peut d'abord remplacer $S$ par le sous-préschéma réduit $S'$ de $S$ ayant $\overline{\{g(x)\}}$ pour espace sous-jacent, $X$ par $g^{-1}(S')$, $Y$ par $h^{-1}(S')$, les fibres de $X$ et $X'$ (resp. $Y$ et $Y'$) aux points de $S'$ étant les mêmes. Autrement dit on peut se borner au cas où $S$ est \emph{intègre} et où $\eta=g(x)=h(y)$ (où $y=f(x)$) est son point générique.

$1^\circ$ Supposons d'abord que $x\in E$. Les anneaux locaux $\sh{O}_{X_\eta,x}$ et $\sh{O}_{Y_\eta,y}$ sont respectivement égaux à $\sh{O}_{X,x}$ et $\sh{O}_{Y,y}$; comme la propriété de lissité d'un morphisme de présentation finie en un point ne dépend que de l'anneau local de ce point et de l'anneau local de son image (17.5.1), on voit que l'hypothèse $x\in E$ revient à dire que le morphisme $f$ est lisse au point $x$; il possède alors encore cette propriété aux points d'un voisinage ouvert de $x$ dans $X$, et il suffit d'appliquer (17.3.3, (iii)) pour obtenir alors la conclusion.

$2^\circ$ Supposons en second lieu que $x\in X-E$, et que le morphisme $f_\eta$ ne soit pas plat au point $x$. La conclusion résulte alors du lemme suivant qui précise (11.2.8):

\begin{lemma}[17.7.11.1]
\label{IV.17.7.11.1}
Soient $g:X\to S$, $h:Y\to S$ deux morphismes localement de présentation finie, $f:X\to Y$ un $S$-morphisme, $\sh{F}$ un $\sh{O}_X$-module quasi-cohérent de présentation finie. Alors l'ensemble $E$ des $x\in X$ tels que $\sh{F}_{g(x)}$ soit $f_{g(x)}$-plat au point $x$ est localement constructible.
\end{lemma}

Raisonnant encore comme au début de (9.9.1) et utilisant (2.5.1), on se ramène au cas où $S$, $X$ et $Y$ sont noethériens; puis on se ramène comme ci-dessus au cas où $S$ est intègre, où $\eta=g(x)=h(f(x))$ est son point générique, et il s'agit de montrer que si $x\in E$ (resp. $x\notin E$) il y a un voisinage $V$ de $x$ dans $\overline{\{x\}}$ contenu dans $E$ (resp. dans $X-E$). Le cas où $x\in E$ est conséquence immédiate de (11.1.1). Pour traiter le cas où $x\notin E$, on raisonne comme dans (9.4.7.1), dont nous conservons les notations, de sorte que l'on peut supposer, en remplaçant éventuellement $Y$ et $X$ par des voisinages de $f(x)$ et $x$ respectivement, qu'il existe deux $\sh{O}_Y$-modules cohérents $\sh{G}$, $\sh{H}$ et un $\sh{O}_Y$-homomorphisme $u:\sh{G}\to\sh{H}$ tels que pour tout $s\in S$, $u_s:\sh{G}_s\to\sh{H}_s$ soit injectif, mais que l'homomorphisme $1\otimes u_\eta:\sh{F}_\eta\otimes_{\sh{O}_{Y_\eta}}\sh{G}_\eta\to\sh{F}_\eta\otimes_{\sh{O}_{Y_\eta}}\sh{H}_\eta$ ne soit pas injectif au point $x$, autrement dit $x\in\operatorname{Supp}(\operatorname{Ker}(1\otimes u_\eta))$; si l'on pose $T=\overline{\{x\}}$, on a donc $\operatorname{Supp}(\operatorname{Ker}(1\otimes u_\eta))\supset T_\eta$. Mais on peut supposer que pour tout $s\in S$, on a $\operatorname{Ker}(1\otimes u_s)=(\operatorname{Ker}(1\otimes u))_s$ (9.4.2), donc

